\chapter[Conclusion]{Conclusion}
\label{Chap:Conclusion}

Ce mémoire étudie la répartition de l'\HI\ cinématiquement étroit et large et ses liens avec la formation stellaire dans 16 galaxies proches. Nous avons décomposé les profils d'\HI\ avec ROHSA, puis comparé les cartes séparées aux cartes de H$_2$ déduites du CO et à celles du SFR. Les composantes cinématiquement froide et chaude sont utilisées comme traceurs indirects du CNM et du WNM.

Les deux composantes atomiques s'étendent sur une grande partie des disques d'\HI\ cartographiés. Sur l'ensemble de l'échantillon, la composante large contribue en médiane à environ 55\% de la masse totale d'\HI, tandis que la composante cinématiquement froide représente les 45\% restants. La composante cinématiquement froide présente un aspect plus filamentaire, mais sa fraction n'est pas systématiquement plus élevée dans les bras spiraux stellaires. Nous ne trouvons pas non plus d'indication claire que l'\HI\ étroit soit généralement plus concentré près de l'émission CO que l'\HI\ large. Cela pourrait s'expliquer par le comportement du gaz atomique pendant la transition atomique-moléculaire, où le gaz atomique reste dans des couches d'écrantage autour des intérieurs moléculaires tandis que du gaz supplémentaire devient moléculaire. L'hydrogène ne contribue en outre plus à l'émission \HI\ une fois converti en H$_2$. De plus, le H$_2$ peut survivre là où le CO est dissocié, de sorte que nos régions faibles en CO ne sont pas nécessairement dépourvues de H$_2$.

Nous avons examiné les relations de Kennicutt--Schmidt des différentes composantes gazeuses. Leur comportement dépend de la population de galaxies. Dans l'échantillon combiné de naines, l'\HI\ étroit est mieux corrélé au SFR que l'\HI\ total ou large. Dans les spirales, l'\HI\ large suit mieux le SFR que l'\HI\ étroit. Cette différence pourrait venir de l'élargissement turbulent des raies du gaz froid près des régions de formation stellaire, qui conduit à en classer une partie dans la composante large. Comme l'ont établi les études antérieures, le gaz moléculaire reste plus étroitement lié au SFR dans la plupart des spirales et présente une relation de Kennicutt--Schmidt approximativement linéaire \citep{Fbigiel2008}.

L'efficacité de formation stellaire des composantes étroite et large de l'\HI\ diminue vers l'extérieur du disque. Dans les naines, celle de l'\HI\ étroit est plus élevée que celle de l'\HI\ large dans tout le disque. Dans les spirales, elle est plus élevée dans le disque interne, puis les efficacités des deux composantes se rapprochent à plus grand rayon. Ce rapprochement pourrait s'expliquer par une moindre contribution de l'élargissement turbulent dans le disque externe, en accord avec la diminution radiale des largeurs de raie d'\HI\ rapportée par \citet{Tamburro2009}. Une part plus importante du gaz froid pourrait alors être affectée à la composante étroite plutôt qu'à la composante large, rapprochant leurs contributions relatives et les SFE mesurées.

Nos données ne permettent pas de localiser la transition \HI--H$_2$, puisque l'hydrogène moléculaire peut être présent sans émission CO détectable. Des observations à meilleures résolutions spatiale et spectrale, des contraintes indépendantes sur les phases thermiques et des mesures du gaz moléculaire tenant compte de la métallicité aideraient à préciser le rôle du gaz atomique froid. Un échantillon plus étendu de galaxies naines permettrait également de déterminer si la plus forte relation combinée de l'\HI\ étroit est une propriété générale de cette population.
